\begin{lemma}[A $C^1$ weak solution is $p$-harmonic with its classical gradient]
Let $p\in\mathbb{R}$ with $p>1$, let $\Omega\subseteq\mathrm{Plane}$, let
$u\colon\mathrm{Plane}\to\mathbb{R}$ and let $g\colon\mathrm{Plane}\to\mathrm{Plane}$.  Assume
\begin{enumerate}
\item $u$ is $p$-harmonic in $\Omega$ with weak gradient $g$;
\item $u$ is of class $C^{1}$ on $\Omega$.
\end{enumerate}
Then
\begin{enumerate}
\item[(a)] $u$ is $p$-harmonic in $\Omega$ with weak gradient equal to its own classical
gradient $Du$; and
\item[(b)] $g(x)=Du(x)$ for Lebesgue-almost every $x\in\Omega$; that is, for almost every
$x\in\mathrm{Plane}$, if $x\in\Omega$ then $g(x)=Du(x)$.
\end{enumerate}
\end{lemma}
\begin{proof}
\medskip\noindent\emph{A remark on real powers, recorded once and used twice below.}  One fact
about real powers is used below which is \emph{not} among the listed clauses of the background
proposition on elementary properties of real powers, so we record it here rather than attaching a
citation the proposition does not support.
\begin{itemize}
\item[(P1)] \emph{Integer exponents.}  For $t\ge0$ and $n\in\mathbb{N}$ the real power $t^{a}$
at $a:=n$ agrees with the $n$-fold product $t^{n}$: for $t>0$ this is part of what it means for
$t^{a}$ to be the real power, which extends the integer power, and for $t=0$ it is the stated
convention, since clause (1) gives $0^{n}=0$ for $n\ge1$ while $0^{0}=1$ is the empty product.
Only the case $n=1$ is used below, namely $t^{1}=t$ for every $t\ge0$.
\end{itemize}

We refer to the six conditions in the definition of a $p$-harmonic function with weak gradient,
\noderef{IsPHarm}, by their numbers (1)--(6), and write $A_p$ for the $p$-flux \noderef{PFlux}.
By condition (1) for the datum of hypothesis (1), the set $\Omega$ is open, hence Lebesgue
measurable by clause (2) of the background proposition on integrability of continuous and
compactly supported functions.  By hypothesis (2) and clause (2) of the background proposition
on locality of $C^{1}$ regularity, $u$ is differentiable at every point of $\Omega$ and both
$u$ and $Du$ are continuous on $\Omega$.

\medskip\noindent\emph{Step 1: integration by parts against a test function.}
\emph{Claim.}  For every $\varphi\colon\mathrm{Plane}\to\mathbb{R}$ of class $C^\infty$ with
compact support and $\operatorname{supp}\varphi\subseteq\Omega$,
\begin{equation}\label{eq:pwgc-ibp}
  \int_\Omega u\,D\varphi\,dx=-\int_\Omega \varphi\,Du\,dx .
\end{equation}

\emph{Proof of the claim.}  Put $K_0:=\operatorname{supp}\varphi$, a compact subset of
$\Omega$, and
\[
  \psi:=u\,\varphi .
\]

\emph{(i) $\psi$ has compact support.}  If $\varphi(x)=0$ then $\psi(x)=0$, so
$\{\psi\neq0\}\subseteq\{\varphi\neq0\}$ and hence
$\operatorname{supp}\psi\subseteq\operatorname{supp}\varphi=K_0$, which is compact.

\emph{(ii) $\psi$ is of class $C^{1}$ on all of $\mathrm{Plane}$.}  The two sets $\Omega$ and
$\mathrm{Plane}\setminus K_0$ are open and cover $\mathrm{Plane}$, because $K_0\subseteq\Omega$.
On $\Omega$ the function $u$ is of class $C^{1}$ by hypothesis (2) and $\varphi$ is of class
$C^\infty$, hence of class $C^{1}$, so $\psi=u\,\varphi$ is of class $C^{1}$ on $\Omega$ by
clause (4) of the background proposition on the product rule for the classical gradient.  On
$\mathrm{Plane}\setminus K_0$ we have $\varphi\equiv0$, since $K_0$ is the closure of
$\{\varphi\neq0\}$, so $\psi\equiv0$ there, and the zero function is of class $C^{1}$ by the
same clause.  Every point of $\mathrm{Plane}$ therefore has an open neighbourhood on which
$\psi$ is of class $C^{1}$, so clause (1) of the locality proposition, applied with the open set
$\mathrm{Plane}$, gives that $\psi$ is of class $C^{1}$ on $\mathrm{Plane}$.

\emph{(iii) The product rule, globally.}  Let $x\in\Omega$.  Both $u$ and $\varphi$ are
differentiable at $x$, so clause (1) of the product-rule proposition gives
\[
  D\psi(x)=\varphi(x)\,Du(x)+u(x)\,D\varphi(x).
\]
Let now $x\in\mathrm{Plane}\setminus K_0$.  On that open set $\psi\equiv0$ and $\varphi\equiv0$,
so clause (3) of the product-rule proposition gives $D\psi(x)=0$ and $D\varphi(x)=0$, while
$\varphi(x)=0$; hence both sides of the display vanish.  Since
$\Omega\cup(\mathrm{Plane}\setminus K_0)=\mathrm{Plane}$,
\begin{equation}\label{eq:pwgc-prod}
  D\psi(x)=\varphi(x)\,Du(x)+u(x)\,D\varphi(x)\qquad\text{for every }x\in\mathrm{Plane}.
\end{equation}

\emph{(iv) The mean of $D\psi$ vanishes.}  By (i) and (ii), $\psi$ is of class $C^{1}$ with
compact support, so \noderef{TestFunctionGradientIntegral} applies to $\psi$ and gives
$\int_{\mathrm{Plane}}D\psi\,dx=0$.  Substituting \eqref{eq:pwgc-prod},
\begin{equation}\label{eq:pwgc-zero}
  \int_{\mathrm{Plane}}\bigl(\varphi\,Du+u\,D\varphi\bigr)\,dx=0 .
\end{equation}

\emph{(v) Localization.}  The functions $\varphi\,Du$ and $u\,D\varphi$ are continuous on
$\Omega$, since $u$ and $Du$ are continuous on $\Omega$ and $\varphi$, $D\varphi$ are continuous
on $\mathrm{Plane}$.  Both vanish off $K_0$: for $x\notin K_0$ we have $\varphi(x)=0$ and, by
clause (3) of the product-rule proposition applied on the open set
$\mathrm{Plane}\setminus K_0$, also $D\varphi(x)=0$.  Each is therefore bounded on the compact
set $K_0\subseteq\Omega$, hence integrable there, by clause (3) of the integrability
proposition; and since each vanishes off $K_0$, clause (4) of that proposition --- applied once
with $S:=\mathrm{Plane}$ and once with $S:=\Omega$, both containing $K_0$ --- gives
\[
  \int_{\mathrm{Plane}}\varphi\,Du=\int_{K_0}\varphi\,Du=\int_\Omega\varphi\,Du,
  \qquad
  \int_{\mathrm{Plane}}u\,D\varphi=\int_{K_0}u\,D\varphi=\int_\Omega u\,D\varphi .
\]
Both summands in \eqref{eq:pwgc-zero} being integrable, the integral splits, and the two
displays give $\int_\Omega u\,D\varphi+\int_\Omega\varphi\,Du=0$.  This is
\eqref{eq:pwgc-ibp}, and the claim is proved.

\medskip\noindent\emph{Step 2: the two weak-gradient identities force $\varphi\,(g-Du)$ to
integrate to zero.}
Let $\varphi$ be as in Step 1, and keep $K_0=\operatorname{supp}\varphi$ and
$B':=\sup_x\lvert\varphi(x)\rvert<\infty$, finite because $\varphi$ is continuous and vanishes
off the compact $K_0$.  Condition (5) of \noderef{IsPHarm} for the datum of hypothesis (1)
gives $\int_\Omega u\,D\varphi=-\int_\Omega\varphi\,g$.  Comparing with \eqref{eq:pwgc-ibp},
\begin{equation}\label{eq:pwgc-same}
  \int_\Omega\varphi\,g\,dx=\int_\Omega\varphi\,Du\,dx .
\end{equation}
Both integrands are integrable on $\Omega$.  For $\varphi\,Du$ this is Step 1(v).  For
$\varphi\,g$: the field $g$ is almost everywhere strongly measurable with respect to Lebesgue
measure restricted to $\Omega$ by condition (3), hence so is $\varphi\,g$; taking $a=1$ in
clause (4) of the background proposition on elementary properties of real powers --- legitimate
since $0\le1\le p$, as $p>1$ --- gives $t^{1}\le1+t^{p}$ for every $t\ge0$, and $t^{1}=t$ by
(P1), so $t\le1+t^{p}$ for every $t\ge0$; applying this at $t:=\lvert\eta\rvert\ge0$ gives
$\lvert\eta\rvert\le1+\lvert\eta\rvert^{p}$ for every vector $\eta$, so
\[
  \lvert\varphi(y)\,g(y)\rvert\le B'\bigl(1+\lvert g(y)\rvert^{p}\bigr),
\]
and the right-hand side is integrable on $K_0$ because $\lvert g\rvert^{p}$ is, by condition
(4), and $\lvert K_0\rvert<\infty$; clause (3) of the integrability proposition then gives
integrability of $\varphi\,g$ on $K_0$, and clause (4) extends it to $\Omega$ and to
$\mathrm{Plane}$, since $\varphi\,g$ vanishes off $K_0$.

We may therefore subtract the two sides of \eqref{eq:pwgc-same}:
\[
  \int_\Omega\varphi(x)\,\bigl(g(x)-Du(x)\bigr)\,dx=0 .
\]
The integrand vanishes off $K_0\subseteq\Omega$, so clause (4) of the integrability proposition
with $S:=\mathrm{Plane}$ gives
\begin{equation}\label{eq:pwgc-test}
  \int_{\mathrm{Plane}}\varphi(x)\,\bigl(g(x)-Du(x)\bigr)\,dx=0
  \qquad\text{for every test function }\varphi\text{ with }\operatorname{supp}\varphi\subseteq\Omega .
\end{equation}

\medskip\noindent\emph{Step 3: part (b), by the fundamental lemma.}
Put $f:=g-Du$, a map $\mathrm{Plane}\to\mathrm{Plane}$.  We check that $f$ is integrable on
every compact subset of $\Omega$, which for the open set $\Omega$ is exactly local
integrability on $\Omega$ as recorded in the fundamental-lemma background proposition.  Let
$K\subseteq\Omega$ be compact.  The field $g$ is almost everywhere strongly measurable on
$\Omega$ by condition (3) and satisfies $\lvert g\rvert\le1+\lvert g\rvert^{p}$ pointwise by
clause (4) of the elementary-powers proposition with $a=1$, legitimate since $0\le1\le p$,
together with $t^{1}=t$ from (P1) as in Step 2; the right-hand side is
integrable on $K$ by condition (4) and $\lvert K\rvert<\infty$, so clause (3) of the
integrability proposition makes $g$ integrable on $K$.  The field $Du$ is continuous on
$\Omega$, hence bounded and integrable on $K$ by clause (3) of the integrability proposition.
So $f=g-Du$ is integrable on $K$.

By \eqref{eq:pwgc-test}, $\int_{\mathrm{Plane}}\varphi\,f\,dx=0$ for every
$\varphi\colon\mathrm{Plane}\to\mathbb{R}$ of class $C^\infty$ with compact support and
$\operatorname{supp}\varphi\subseteq\Omega$.  The fundamental lemma of the calculus of
variations, applied with $U:=\Omega$, $F:=\mathrm{Plane}$ and this $f$, therefore gives
$f(x)=0$ for almost every $x\in\Omega$, that is
\begin{equation}\label{eq:pwgc-ae}
  g(x)=Du(x)\qquad\text{for almost every }x\in\Omega .
\end{equation}
This is part (b).

\medskip\noindent\emph{Step 4: part (a).}
We verify the six conditions of \noderef{IsPHarm} for $u$ and $Du$ on $\Omega$.

(1) $\Omega$ is open, by condition (1) for the datum of hypothesis (1); this condition does not
involve the gradient.

(2) $u$ is continuous on $\Omega$, by condition (2) for that datum.

(3) By hypothesis (2) and clause (2) of the locality proposition, $Du$ is continuous on
$\Omega$; and $\Omega$ is measurable.  Clause (3) of the integrability proposition then gives
that $Du$ is almost everywhere strongly measurable with respect to Lebesgue measure restricted
to $\Omega$.

(4) Let $K\subseteq\Omega$ be compact.  Since $Du$ is continuous on $\Omega$, clause (3) of the
integrability proposition gives a bound $\lvert Du(y)\rvert\le B$ for $y\in K$, and we may
assume $B\ge1$.  As $p>0$, clause (3) of the elementary-powers proposition (monotonicity of
$t\mapsto t^{p}$ on $[0,\infty)$) gives $\lvert Du(y)\rvert^{p}\le B^{p}$ for $y\in K$.  The
function $y\mapsto\lvert Du(y)\rvert^{p}$ is measurable on $K$, being a composition of the
continuous maps $\eta\mapsto\lvert\eta\rvert$ and $t\mapsto t^{p}$ with $Du$, and it is
dominated by the constant $B^{p}$, which is integrable on $K$ because
$\lvert K\rvert<\infty$.  Clause (3) of the integrability proposition therefore makes
$\lvert Du\rvert^{p}$ integrable on $K$.

(5) This is exactly \eqref{eq:pwgc-ibp} of Step 1.

(6) Let $\varphi$ be of class $C^\infty$ with compact support and
$\operatorname{supp}\varphi\subseteq\Omega$.  By \eqref{eq:pwgc-ae} we have $g=Du$ almost
everywhere on $\Omega$, so by the composition remark in clause (6) of the background
proposition on restriction and congruence, $A_p\circ g=A_p\circ Du$ almost everywhere on
$\Omega$, and hence
$\bigl\langle A_p(g),D\varphi\bigr\rangle=\bigl\langle A_p(Du),D\varphi\bigr\rangle$ almost
everywhere on $\Omega$.  The same clause (6), which says that almost everywhere equal functions
have the same integral, gives
\[
  \int_\Omega\bigl\langle A_p(Du(x)),D\varphi(x)\bigr\rangle\,dx
  =\int_\Omega\bigl\langle A_p(g(x)),D\varphi(x)\bigr\rangle\,dx=0 ,
\]
the last equality being condition (6) for the datum of hypothesis (1).

\medskip\noindent All six conditions hold, so $u$ is $p$-harmonic in $\Omega$ with weak gradient
$Du$.  This is part (a), and part (b) was proved in Step 3.
\end{proof}
